\subsection{Equações da cinemática}
A atitude de uma espaçonave pode ser descrita pelos ângulos de Euler e pelos quaternions. 
Cada representação possui vantagens e limitações quanto aos parâmetros geométricos, à presença de singularidades e ao custo computacional, o que impacta as simulações numéricas.
% Nesta subseção do trabalho serão apresentadas as principais formas de representação e suas respectivas relações.
% Será dado ênfase nos quaternions unitários, pois são utilizados em aplicações aeroespaciais devido à eficiência e relativo baixo custo computacional, quando comparado com os ângulos de Euler.

%==================================================
%----- Ang de Euler
%==================================================
\subsubsection{Ângulos de Euler}
Neste trabalho, foi adotado a sequência de rotações extrínsecas ZYX, conhecida como \emph{yaw, pitch e roll}.
% A rotação em torno do eixo $Z$ ($\psi$) é aplicada primeiro, seguida pela rotação em torno do eixo $Y$ ($\theta$), e por fim pela rotação em torno do eixo $X$ ($\phi$).
Embora a interpretação ser direta (roll, pitch e yaw), os ângulos de Euler podem sofrer com a ocorrência de singularidades (\emph{gimbal lock})
\footnote{A singularidade ocorre quando $\theta = \pm 90^{\circ}$, fazendo com que dois eixos de rotação se tornem linearmente dependentes, ou seja, é perdido um grau de liberdade.}.

% Além dos ângulos de Euler e dos quaternions, outra forma equivalente de representar a orientação é por meio das matrizes de rotação. Essas matrizes são utilizadas em robótica e em transformações entre os sistemas de referência, por exemplo do referencial inercial para o referencial LVLH. As formulações completas, são apresentadas no Apêndice \ref{ApendiceE}.

%==================================================
%----- Quaternions
%==================================================
\subsubsection{Quaternions}
\label{sec:Quaternions}
Os quaternions
% são empregados em aplicações aeroespaciais por causa de sua eficiência computacional e por evitarem
evitam as singularidades associadas aos ângulos de Euler.
Um quaternion é definido como:
\begin{equation}
    \label{eq:quaternion_unitario_definicao}
\mathbf{q}
=
\begin{bmatrix}
q_0 \\ q_1 \\ q_2 \\ q_3
\end{bmatrix}
=
\begin{bmatrix}
\cos\frac{\alpha}{2} \\
u_x \sin\frac{\alpha}{2} \\
u_y \sin\frac{\alpha}{2} \\
u_z \sin\frac{\alpha}{2}
\end{bmatrix},
\end{equation}
onde $\alpha$ é o ângulo de rotação e $\vec{u} = (u_x,u_y,u_z)^T$ é o eixo unitário em torno do qual a rotação ocorre.
É importante esclarecer que
$\mathbf{q}\in\mathbb{H}$ com $|\mathbf{q}|=1$.

\subsubsection{Produto de quaternions}
Dados dois quaternions
$p = [p_0,\vec{p}_v]^T$ e $q = [q_0,\vec{q}_v]^T$, 
o produto $p \otimes q$ é definido como:
\begin{equation}
p \otimes q
=
\begin{bmatrix}
p_0 q_0 - \vec{p_v}^{T}\vec{q}_v \\
p_0 \vec{q}_v + q_0 \vec{p}_v + \vec{p}_v \times \vec{q}_v
\end{bmatrix}.
\label{eq:quat_produto}
\end{equation}

%==================================================
%----- Rel. quaternion e matriz de rotação
%==================================================
\subsubsection{Relação entre quaternion e matriz de rotação}
Um quaternion unitário é definido como:
\begin{equation}
q = q_0
+ q_1 \mathbf{i}
+ q_2\mathbf{j}
+ q_3\mathbf{k},
\label{eq:def_quaternion}
\end{equation}
onde $q_0$ é a parte escalar e $\vec{q}_v = [q_1; q_2; q_3]^T$ é a parte vetorial.  

A rotação de um vetor $\vec{v}$ pode ser escrita como:
\begin{equation}
\vec{v}' = (q_0^2 -
\vec q_v{^T} \vec{q}_v)\,\vec{v} 
+ 2(\vec{q}_v \cdot \vec{v})\,\vec{q}_v
+ 2q_0 (\vec{q}_v \times \vec{v}).
\label{eq:rodrigues_quaternion}
\end{equation}

Reescrevendo \eqref{eq:rodrigues_quaternion} em forma matricial, obtém-se:

\begin{equation}
\mathbf{R}(q) =
(q_0^2 - \vec{q}_v{^T} \vec{q}_v)\mathbf{I}_3 
+ 2 \vec{q}_v \vec{q}_v{^T}
+ 2q_0 [\vec{q}_v]_\times,
\label{eq:rot_matrix_general}
\end{equation}

onde $[\vec{q}_v]_\times$ é a matriz antissimétrica
% \fn{Essa matriz auxilia nas transformações de coordenadas e na formulação das equações dinâmicas, sendo definida como uma matriz associada a um vetor qualquer $\vec a = [x, y, z]^T$.}
($\hat{\vec{q}}_v$), associada ao produto vetorial:

\begin{equation}
\hat{\vec{q}}_v =
[\vec{q}_v]_\times =
\begin{bmatrix}
0 & -q_3 & q_2 \\
q_3 & 0 & -q_1 \\
-q_2 & q_1 & 0
\end{bmatrix}.
\label{eq:skew_matrix}
\end{equation}

Expandindo os termos de \eqref{eq:rot_matrix_general}, obtém-se:
\begin{equation}
\mathbf{R}(q) =
\begin{bmatrix}
1 - 2(q_2^2+q_3^2) & 2(q_1q_2 - q_0q_3) & 2(q_1q_3+q_0q_2) \\
2(q_1q_2+q_0q_3)   & 1 - 2(q_1^2+q_3^2) & 2(q_2q_3 - q_0q_1) \\
2(q_1q_3 - q_0q_2) & 2(q_2q_3+q_0q_1)   & 1 - 2(q_1^2+q_2^2)
\end{bmatrix}.
\label{eq:rot_matrix_expandida}
\end{equation}

Assim, \eqref{eq:rot_matrix_expandida} mostra a forma expandida da matriz de rotação a partir de um quaternion unitário.

\subsubsection{Matriz antissimétrica} 
A matriz antissimétrica (\emph{skew-symmetric matrix}), apresentada em \eqref{eq:skew_matrix}, auxilia nas transformações de coordenadas e na formulação das equações dinâmicas.  
Ela é definida a partir de um vetor genérico
$\vec{a} = [x,y,z]^T$, como:

\begin{equation}
    \label{eq:matriz_assimetrica}
    \hat{\vec{a}} = 
    \begin{bmatrix}
        0   & -z  &  y \\
        z   &  0  & -x \\
        -y  &  x  &  0
    \end{bmatrix}.
\end{equation}

Essa matriz possui a propriedade fundamental de que
$\hat{\vec{a}}\vec{b} = \vec{a} \times \vec{b}$,
ou seja, a multiplicação matricial equivale ao produto vetorial tradicional.

%==================================================
%----- Cinemática diferencial
%==================================================
\subsection{Cinemática diferencial dos quaternions}
\label{sec:cinematica_diferencial_quaternions}

Considerando a velocidade angular $\vec \omega$ expressa na base ($\{0\}$), é possível determinar a matriz de atitude como:

\begin{equation}
    \label{eq:dotR_base}
\dot{\mathbf{R}}
=
\mathbf{R}\,[\vec{\omega}]_{\times},
\end{equation}

onde $[\vec{\omega}]_{\times}$ denota a matriz antissimétrica \eqref{eq:skew_matrix}.
Por sua vez, a cinemática do quaternion unitário
$\vec{q}=[q_{0},\vec{\epsilon}\,]^{T}\,$
é dada por:

% \begin{equation}
%     \label{eq:definicao_qdot}
% \dot{\vec{q}}
% =
% \tfrac{1}{2}\,\Omega(\vec{q})\,\vec{\omega},
% \qquad
% \dot{\vec{q}}
% =
% \tfrac{1}{2}\,\mathbf{Q}(\vec{\omega})\,\vec{q},
% \end{equation}

\begin{equation}
    \label{eq:definicao_qdot}
\dot{\vec{q}}
=
\tfrac{1}{2}\,\Omega(\vec{q})\,\vec{\omega},
\end{equation}

% onde a matriz de propagação $\mathbf{Q}$ dos quaternions é dada como:
% \begin{equation}
% \mathbf{Q}(\vec{\omega}) =
% \begin{bmatrix}
% 0 & -\vec{\omega}^T \\
% \vec{\omega} & -[\vec{\omega}]_\times
% \end{bmatrix},
% \end{equation}

% e sendo expandida como:
% \begin{equation}
% \mathbf{Q}(\vec{\omega}) =
% \begin{bmatrix}
% 0 & -\omega_x & -\omega_y & -\omega_z \\
% \omega_x & 0 & \omega_z & -\omega_y \\
% \omega_y & -\omega_z & 0 & \omega_x \\
% \omega_z & \omega_y & -\omega_x & 0
% \end{bmatrix}.
% \end{equation}

%%%%% Como a velocidade angular $\vec \omega$ possui dimensão igual a 3, deve ser transformada para 4 dimensões, sendo expressa matematicamente como:
sendo a matriz de mapeamento
$\Omega(\vec{q})\in\mathbb{R}^{4\times 3}\,$
que converte a velocidade angular
$\vec{\omega}\in\mathbb{R}^{3}$
em
$\dot{\vec{q}}\in\mathbb{R}^{4}$:

\begin{equation}
    \label{eq:Omega_def}
\Omega(\vec{q})
=
\begin{bmatrix}
-\vec{\epsilon\,}^{T} \\[2pt]
q_0\,\mathbf{I}_3
+ [\vec{\epsilon}]_{\times}
\end{bmatrix}.
\end{equation}

% Quando a velocidade angular é expressa na referência inercial, usa-se a forma equivalente, que é consistente com a escolha do sistema de referência.
% \begin{equation}
%     \label{eq:equivalente_velocAngular}
% \dot{\mathbf{R}}=[\vec{\omega}]_{\times}\mathbf{R}
% \end{equation}

Por fim, para obter estabilidade numérica, é realizada a normalização periódica de $q$, pois eventuais erros de integração podem violar a unidade da norma.

%==================================================
%----- Normalização e consistência numérica
%==================================================
\subsubsection{Normalização}
A principal restrição dos quaternions é que eles devem permanecer normalizados:

\begin{equation}
\label{eq:normalizacao_quaternion}
\lVert \mathbf{q} \rVert^2 =
q_0^2
+ q_1^2
+ q_2^2
+ q_3^2
= 1.
\end{equation}

Durante as simulações numéricas, erros de arredondamento e de integração numérica podem violar essa condição.
Além disso, a normalização não altera a atitude, apenas mitiga inconsistências numéricas, e para preservá-las, adota-se a normalização periódica do vetor $\mathbf{q}$ ao longo das simulações:

\begin{equation}
\mathbf{q}
\leftarrow
\frac{\mathbf{q}}{\lVert \mathbf{q} \rVert}.
\end{equation}

Nesta dissertação será adotada a convenção de notação do quaternion como
\begin{equation}
q = [\,w,\,x,\,y,\,z\,]^T,
\label{eq:quat_convencao}
\end{equation}
onde $w$ representa a parte escalar e $(x,y,z)$ correspondem à parte vetorial.

% \subsection{Vantagens do uso de quaternions}
% Dentre as vantagens principais destacam-se: a ausência de singularidades como o \emph{gimbal lock}, a eficiência computacional, a representação compacta (quatro parâmetros em vez de nove da matriz de rotação) e a estabilidade numérica em simulações de longa duração.

% Nesta seção foi apresentado o conjunto de equações que descreve a cinemática da atitude, ou seja, a evolução temporal da orientação em função da velocidade angular.
% A seguir, serão apresentadas as equações dinâmicas de Euler, que relacionam as acelerações angulares aos torques aplicados na espaçonave.